% Options for packages loaded elsewhere
% Options for packages loaded elsewhere
\PassOptionsToPackage{unicode}{hyperref}
\PassOptionsToPackage{hyphens}{url}
\PassOptionsToPackage{dvipsnames,svgnames,x11names}{xcolor}
%
\documentclass[
  12pt,
  letterpaper,
  DIV=11,
  numbers=noendperiod]{scrartcl}
\usepackage{xcolor}
\usepackage{amsmath,amssymb}
\setcounter{secnumdepth}{5}
\usepackage{iftex}
\ifPDFTeX
  \usepackage[T1]{fontenc}
  \usepackage[utf8]{inputenc}
  \usepackage{textcomp} % provide euro and other symbols
\else % if luatex or xetex
  \usepackage{unicode-math} % this also loads fontspec
  \defaultfontfeatures{Scale=MatchLowercase}
  \defaultfontfeatures[\rmfamily]{Ligatures=TeX,Scale=1}
\fi
\usepackage{lmodern}
\ifPDFTeX\else
  % xetex/luatex font selection
  \setmainfont[]{Times New Roman}
\fi
% Use upquote if available, for straight quotes in verbatim environments
\IfFileExists{upquote.sty}{\usepackage{upquote}}{}
\IfFileExists{microtype.sty}{% use microtype if available
  \usepackage[]{microtype}
  \UseMicrotypeSet[protrusion]{basicmath} % disable protrusion for tt fonts
}{}
\makeatletter
\@ifundefined{KOMAClassName}{% if non-KOMA class
  \IfFileExists{parskip.sty}{%
    \usepackage{parskip}
  }{% else
    \setlength{\parindent}{0pt}
    \setlength{\parskip}{6pt plus 2pt minus 1pt}}
}{% if KOMA class
  \KOMAoptions{parskip=half}}
\makeatother
% Make \paragraph and \subparagraph free-standing
\makeatletter
\ifx\paragraph\undefined\else
  \let\oldparagraph\paragraph
  \renewcommand{\paragraph}{
    \@ifstar
      \xxxParagraphStar
      \xxxParagraphNoStar
  }
  \newcommand{\xxxParagraphStar}[1]{\oldparagraph*{#1}\mbox{}}
  \newcommand{\xxxParagraphNoStar}[1]{\oldparagraph{#1}\mbox{}}
\fi
\ifx\subparagraph\undefined\else
  \let\oldsubparagraph\subparagraph
  \renewcommand{\subparagraph}{
    \@ifstar
      \xxxSubParagraphStar
      \xxxSubParagraphNoStar
  }
  \newcommand{\xxxSubParagraphStar}[1]{\oldsubparagraph*{#1}\mbox{}}
  \newcommand{\xxxSubParagraphNoStar}[1]{\oldsubparagraph{#1}\mbox{}}
\fi
\makeatother


\usepackage{longtable,booktabs,array}
\newcounter{none} % for unnumbered tables
\usepackage{calc} % for calculating minipage widths
% Correct order of tables after \paragraph or \subparagraph
\usepackage{etoolbox}
\makeatletter
\patchcmd\longtable{\par}{\if@noskipsec\mbox{}\fi\par}{}{}
\makeatother
% Allow footnotes in longtable head/foot
\IfFileExists{footnotehyper.sty}{\usepackage{footnotehyper}}{\usepackage{footnote}}
\makesavenoteenv{longtable}
\usepackage{graphicx}
\makeatletter
\newsavebox\pandoc@box
\newcommand*\pandocbounded[1]{% scales image to fit in text height/width
  \sbox\pandoc@box{#1}%
  \Gscale@div\@tempa{\textheight}{\dimexpr\ht\pandoc@box+\dp\pandoc@box\relax}%
  \Gscale@div\@tempb{\linewidth}{\wd\pandoc@box}%
  \ifdim\@tempb\p@<\@tempa\p@\let\@tempa\@tempb\fi% select the smaller of both
  \ifdim\@tempa\p@<\p@\scalebox{\@tempa}{\usebox\pandoc@box}%
  \else\usebox{\pandoc@box}%
  \fi%
}
% Set default figure placement to htbp
\def\fps@figure{htbp}
\makeatother


% definitions for citeproc citations
\NewDocumentCommand\citeproctext{}{}
\NewDocumentCommand\citeproc{mm}{%
  \begingroup\def\citeproctext{#2}\cite{#1}\endgroup}
\makeatletter
 % allow citations to break across lines
 \let\@cite@ofmt\@firstofone
 % avoid brackets around text for \cite:
 \def\@biblabel#1{}
 \def\@cite#1#2{{#1\if@tempswa , #2\fi}}
\makeatother
\newlength{\cslhangindent}
\setlength{\cslhangindent}{1.5em}
\newlength{\csllabelwidth}
\setlength{\csllabelwidth}{3em}
\newenvironment{CSLReferences}[2] % #1 hanging-indent, #2 entry-spacing
 {\begin{list}{}{%
  \setlength{\itemindent}{0pt}
  \setlength{\leftmargin}{0pt}
  \setlength{\parsep}{0pt}
  % turn on hanging indent if param 1 is 1
  \ifodd #1
   \setlength{\leftmargin}{\cslhangindent}
   \setlength{\itemindent}{-1\cslhangindent}
  \fi
  % set entry spacing
  \setlength{\itemsep}{#2\baselineskip}}}
 {\end{list}}
\usepackage{calc}
\newcommand{\CSLBlock}[1]{\hfill\break\parbox[t]{\linewidth}{\strut\ignorespaces#1\strut}}
\newcommand{\CSLLeftMargin}[1]{\parbox[t]{\csllabelwidth}{\strut#1\strut}}
\newcommand{\CSLRightInline}[1]{\parbox[t]{\linewidth - \csllabelwidth}{\strut#1\strut}}
\newcommand{\CSLIndent}[1]{\hspace{\cslhangindent}#1}



\setlength{\emergencystretch}{3em} % prevent overfull lines

\providecommand{\tightlist}{%
  \setlength{\itemsep}{0pt}\setlength{\parskip}{0pt}}



 


\usepackage{tcolorbox}
\usepackage{amssymb}
\usepackage{yfonts}
\usepackage{bm}


\newtcolorbox{greybox}{
  colback=white,
  colframe=blue,
  coltext=black,
  boxsep=5pt,
  arc=4pt}
  
\newcommand{\sectionbreak}{\clearpage}

 
\newcommand{\ds}[4]{\sum_{{#1}=1}^{#3}\sum_{{#2}=1}^{#4}}
\newcommand{\us}[3]{\mathop{\sum\sum}_{1\leq{#2}<{#1}\leq{#3}}}

\newcommand{\ol}[1]{\overline{#1}}
\newcommand{\ul}[1]{\underline{#1}}

\newcommand{\amin}[1]{\mathop{\text{argmin}}_{#1}}
\newcommand{\amax}[1]{\mathop{\text{argmax}}_{#1}}

\newcommand{\ci}{\perp\!\!\!\perp}

\newcommand{\mc}[1]{\mathcal{#1}}
\newcommand{\mb}[1]{\mathbb{#1}}
\newcommand{\mf}[1]{\mathfrak{#1}}

\newcommand{\eps}{\epsilon}
\newcommand{\lbd}{\lambda}
\newcommand{\alp}{\alpha}
\newcommand{\df}{=:}
\newcommand{\am}[1]{\mathop{\text{argmin}}_{#1}}
\newcommand{\ls}[2]{\mathop{\sum\sum}_{#1}^{#2}}
\newcommand{\ijs}{\mathop{\sum\sum}_{1\leq i<j\leq n}}
\newcommand{\jis}{\mathop{\sum\sum}_{1\leq j<i\leq n}}
\newcommand{\sij}{\sum_{i=1}^n\sum_{j=1}^n}
	
\KOMAoption{captions}{tableheading}
\makeatletter
\@ifpackageloaded{caption}{}{\usepackage{caption}}
\AtBeginDocument{%
\ifdefined\contentsname
  \renewcommand*\contentsname{Table of contents}
\else
  \newcommand\contentsname{Table of contents}
\fi
\ifdefined\listfigurename
  \renewcommand*\listfigurename{List of Figures}
\else
  \newcommand\listfigurename{List of Figures}
\fi
\ifdefined\listtablename
  \renewcommand*\listtablename{List of Tables}
\else
  \newcommand\listtablename{List of Tables}
\fi
\ifdefined\figurename
  \renewcommand*\figurename{Figure}
\else
  \newcommand\figurename{Figure}
\fi
\ifdefined\tablename
  \renewcommand*\tablename{Table}
\else
  \newcommand\tablename{Table}
\fi
}
\@ifpackageloaded{float}{}{\usepackage{float}}
\floatstyle{ruled}
\@ifundefined{c@chapter}{\newfloat{codelisting}{h}{lop}}{\newfloat{codelisting}{h}{lop}[chapter]}
\floatname{codelisting}{Listing}
\newcommand*\listoflistings{\listof{codelisting}{List of Listings}}
\makeatother
\makeatletter
\makeatother
\makeatletter
\@ifpackageloaded{caption}{}{\usepackage{caption}}
\@ifpackageloaded{subcaption}{}{\usepackage{subcaption}}
\makeatother
\usepackage{bookmark}
\IfFileExists{xurl.sty}{\usepackage{xurl}}{} % add URL line breaks if available
\urlstyle{same}
% fallback for those not using the hyperref driver hyperxmp:
\makeatletter
\@ifundefined{xmpquote}{\newcommand{\xmpquote}[1]{#1}}{}
\makeatother
\hypersetup{
  pdftitle={Metric/Nonmetric Elastic MDS},
  pdfauthor={Jan de Leeuw},
  colorlinks=true,
  linkcolor={blue},
  filecolor={Maroon},
  citecolor={Blue},
  urlcolor={Blue},
  pdfcreator={LaTeX via pandoc}}


\title{Metric/Nonmetric Elastic MDS}
\author{Jan de Leeuw}
\date{October 28, 2026}
\begin{document}
\maketitle
\begin{abstract}
We present R and C implementations of the majorization method for metric
and non-metric (ordinal) versions of Elastic MDS, the multidimensional
scaling technique proposed by Victor E. McGee in 1966. The non-metric
version requires a special purpose monotone regression technique. The R
and C versions are compared for speed, with the C version anywhere from
15 to 100 times as fast as the R version.
\end{abstract}

\renewcommand*\contentsname{Table of contents}
{
\setcounter{tocdepth}{3}
\tableofcontents
}

\textbf{Note:} This is a working manuscript which will be
expanded/updated frequently. All suggestions for improvement are
welcome. All Rmd, tex, html, pdf, R, and C files are in the public
domain. Attribution will be appreciated, but is not required. The files
can be found at \url{https://github.com/deleeuw/elastic}

\section{Introduction}\label{introduction}

Early in the history of computerized Multidimensional Scaling (MDS),
between the seminal contributions of Shepard/Kruskal and
Guttman/Lingoes/Roskam, there was the ``elastic method'' proposed by
Victor E. McGee in a series of papers (McGee
(\citeproc{ref-mcgee_65}{1965}), McGee (\citeproc{ref-mcgee_66}{1966}),
McGee (\citeproc{ref-mcgee_67}{1967}), McGee
(\citeproc{ref-mcgee_68}{1968})). The method has been largely forgotten,
but it is worth preserving, because it is different in some important
aspects from the more well-known methods.

The least squares loss function in most MDS methods (Borg and Groenen
(\citeproc{ref-borg_groenen_05}{2005})) can be written as
\begin{equation}
  \sigma(X,\Delta) := \frac{\mathop{\sum\sum}_{1 \leq i<j \leq n}w_{ij}(\delta_{ij} - d_{ij}(X))^2}{\mathop{\sum\sum}_{1 \leq i<j \leq n}w_{ij}\delta_{ij}^2},\label{eq-kruskal}
\end{equation} where \(X\) is the \emph{configuration} of \(n\)
\emph{points} in \(p\) \emph{dimensions}, \(d_{ij}(X)\) is the
\emph{Euclidean distance} between points \(i\) and \(j\) in
configuration \(X\), \(\delta_{ij}\) is the \emph{dissimilarity} between
points \(i\) and \(j\), and \(w_{ij}\) is a \emph{weight}. Weights are
non-negative.

In the \emph{metric} version of the MDS method the dissimilarities are
observed and fixed, and minimization is over configurations only. The
denominator in \eqref{eq-kruskal} is irrelevant for the minimization
problem, but it normalizes the problem in the sense that the minimum of
stress is between zero and one. In the \emph{nonmetric} version
minimization is over both configurations and dissimilarities, with the
constraint that the dissimilarities are monotonic with an observed set
of dissimilarities.

It is true that in the original Kruskal formulation the denominator is
the sum of the squared distances instead. But {de Leeuw}
(\citeproc{ref-deleeuw_U_75a}{1975}) shows that normalizing by using
either the sum of squared distances or the sum of squared
dissimilarities leads to the same solution (up to a scalar
proportionality factor). The smacof program (De Leeuw and Mair
(\citeproc{ref-deleeuw_mair_A_09c}{2009}), {Mair et al.}
(\citeproc{ref-mair_groenen_deleeuw_A_22}{2022})) minimizes the
numerator of \eqref{eq-kruskal} over configurations and dissimilarities,
with the additional constraint that the sum of squares of the
dissimilarities is equal to a constant. Again, this \emph{explicit
normalization} gives the same solution, up to proportionality, as the
original Kruskal formulation that uses \emph{implicit normalization} by
dividing by the sum of squared distances.

In McGee's elastic method the loss is constructed to satisfy two
requirements. The first one is

\begin{quote}
Psychological judgments which indicated relatively great separation of
stimuli should be allowed greater error than judgments indicating close
proximity (l.c. p.~182)
\end{quote}

And the second requirement is the basic MDS requirement that
dimensionality of \(X\) must as low as possible, while still providing a
good fit.

\begin{quote}
A criterion which suggested itself in response to the first requirement
was one based on the physical work done on an elastic spring to stretch
or compress it from an initial length \(\delta_{ij}\) to a final length
\(d_{ij}\). (l.c. p.~183)
\end{quote}

This leads to the loss function \begin{equation}
  \sigma(X,\Delta) := \mathop{\sum\sum}_{1 \leq i<j \leq n}w_{ij}\frac{(\delta_{ij} - d_{ij}(X))^2}{\delta_{ij}^2},\label{eq-mcgee}
\end{equation} which McGee calls \emph{work}. We shall just call it
\emph{stress} (and \emph{elastic stress} or \emph{McGee stress}), using
the more familiar MDS name for loss. We will also continue to use the
symbol \(\sigma\) for any least squares MDS loss function.

The elastic MDS problem is minimization of stress over both \(X\) and
\(\Delta\), where \(\Delta\) must be monotone with the given
dissimilarities. In \eqref{eq-mcgee} the weight \(w_{ij}\) is
interpreted as the modulus of elasticity of the spring \((i,j)\).

In McGee (\citeproc{ref-mcgee_67}{1967}) the alternative loss function
\begin{equation}
  \sigma(X,\Delta) := \mathop{\sum\sum}_{1 \leq i<j \leq n}w_{ij}\frac{(\delta_{ij} - d_{ij}(X))^2}{d_{ij}^2(X)},\label{eq-mcgee2}
\end{equation} is proposed. Minimizing \eqref{eq-mcgee2} seems more
complicated, and we will postpone studying algorithms to minimize it.
Thus we will work with \eqref{eq-mcgee} in this paper.

In McGee's papers the actual algorithm and its implementation are not
described in sufficient detail. Part of the problem is that he is
dealing exclusively with the nonmetric case, in which minimization over
both \(X\) and \(\Delta\) is necessary. In the metric case minimization
is over \(X\) only, and the algorithm is much simpler.

\section{Properties}\label{properties}

If we compare \eqref{eq-kruskal} and \eqref{eq-mcgee} we see one
important difference. The normalization in \eqref{eq-kruskal} is by the
sum of squared dissimilarities, while the normalization in
\eqref{eq-mcgee} is term-by-term, by the squared dissimilarities
themselves. An alternative way of writing the elastic loss function
makes this more clear. \begin{equation}
  \sigma(X,\Delta) := \mathop{\sum\sum}_{1 \leq i<j \leq n}w_{ij}\left( 1 - \frac{d_{ij}(X)}{\delta_{ij}} \right)^2\label{eq-mcgee3}
\end{equation} Thus we see that instead of minimizing the difference of
dissimilarities and distances from zero the elastic method minimizes the
deviations of their ratios from one.

If \(\delta_{ij}\) is multiplied by a positive constant, then so is the
optimum configuration \(X\) and the \(D(X)\). As a consequence, the
minimum of elastic stress does not change, i.e.~is homogeneous of degree
zero in the dissimilarities.

It is clear from \eqref{eq-mcgee3} that the loss function is undefined
if one or more of the dissimilarities are zero. If one of the distances
is zero, then the corresponding term in the loss function is equal to
the corresponding weight, irrespective of what the corresponding
dissimilarity is. Thus the minimum of elastic stress is always less than
or equal to the sum of the weights, its value at \(X = 0\).

If \(\delta_{ij}\) and \(d_{ij}(X)\) are non-zero and close then a first
order Taylor series expansion gives \begin{equation}
  \log d_{ij}(X) - \log\delta_{ij} \approx \frac{1}{\delta_{ij}}(d_{ij}(X) - \delta_{ij}),\label{eq-ddelta}
\end{equation} from which it follows that \begin{equation}
  \sigma(X,\Delta) \approx \mathop{\sum\sum}_{1 \leq i<j \leq n}w_{ij}(\log\delta_{ij} - \log d_{ij}(X))^2.\label{eq-approx}
\end{equation} If the fit is good the elastic stress will be
approximately equal to the logarithmic stress from Ramsay
(\citeproc{ref-ramsay_77}{1977}). Or, to put it differently, minimizing
elastic stress can serve as an approximation to minimizing logarithmic
stress.

\section{Algorithm}\label{algorithm}

In this paper we will give an iterative algorithm for minimizing loss
from \eqref{eq-mcgee} in the metric and non-metric (ordinal) case. A
majorization algorithm for the metric case is already available in the
smacofx package ({Rusch et al.}
(\citeproc{ref-rusch_deleeuw_chen_mair_M_25}{2025}), {Rusch et al.}
(\citeproc{ref-rusch_deleeuw_mair_hornik_26}{2026})). Our new non-metric
algorithm is in the \emph{alternating least squares} family, which means
we alternate minimizing over \(X\) for fixed \(\Delta\) and over
\(\Delta\) for fixed \(X\).

We will start our iterations with the classical metric MDS solution
(Torgerson (\citeproc{ref-torgerson_58}{1958})) for \(\Delta\). We
actually scale that solution by minimizing \begin{equation}
  \sigma(\lambda) := \sum_{k=1}^m \frac{w_k}{\delta_k^2}(\delta_k - \lambda d_k(X))^2\label{eq-lstress}
\end{equation} over \(\lambda\). The minimum is attained at
\begin{equation}
  \hat\lambda := \frac{\sum_{k=1}^m \frac{w_k}{\delta_k}d_k(X)}{\sum_{k=1}^m \frac{w_k}{\delta_k^2}d_k^2(X)}.\label{eq-lbd}
\end{equation}

In iteration \(k\) we perform one majorization step to replace
\(X^{(k)}\) by \(X^{(k+1)}\) to improve loss for fixed \(\Delta^{(k)}\)
and one monotone regression step to replace \(\Delta^{(k)}\) by
\(\Delta^{(k+1)}\) for fixed \(X^{(k+1)}\).

To minimize \eqref{eq-mcgee} over \(X\) for fixed \(\Delta\) we rewrite
loss as \begin{equation}
  \sigma(X,\Delta) = \mathop{\sum\sum}_{1 \leq i<j \leq n}\frac{w_{ij}}{\delta_{ij}^2}(\delta_{ij} - d_{ij}(X))^2.\label{eq-mcgee4}
\end{equation} This can be minimized (or decreased) by using the
standard smacof majorization step with the weights
\(w_{ij}/\delta_{ij}^2\).

To minimize over \(\Delta\) for given \(X\) we define
\(\gamma_{ij} := -\delta_{ij}^{-1}\) and
\(c_{ij}(X) := -d_{ij}^{-1}(X)\). Rewrite loss as \begin{equation}
  \sigma(X,\Gamma) = \mathop{\sum\sum}_{1 \leq i<j \leq n}w_{ij}d_{ij}^2(X)(\gamma_{ij} - c_{ij}(X))^2,\label{eq-mcgee5}
\end{equation} which we must minimize over increasing \(\gamma_{ij}\).
This is just monotone regression with the \(c_{ij}(X)\) as the targets
and with weights \(w_{ij}d_{ij}^2(X)\). After we have found the optimal
\(\hat\gamma_{ij}\) we transform back to
\(\hat\delta_{ij} = -\hat\gamma_{ij}^{-1}\).

\section{Software}\label{software}

The github repository contains R and C versions of the smacofSSElastic
program. Both smacofSSElasticC() and smacofSSElasticR() have the same
default values of the parameters and the two programs are structured the
same way. We iterate a maximum of 1000 iterations until the change in
stress from one iteration to the next is less than \ensuremath{10^{-6}}.
A single majorization update is alternated with a single monotone
regression.

smacofSSElasticC() has a driver in R that sets up the MDS data structure
({de Leeuw} (\citeproc{ref-deleeuw_E_25b}{2025})). Both programsuse a
number of R utilities for data manipulation and to compute the initial
configuration. One important difference is that the R version uses
gpava(), which is written in R, for monotone regression (De Leeuw et al.
(\citeproc{ref-deleeuw_hornik_mair_A_09}{2009})), while the R version
uses the C routine monotone() from Busing
(\citeproc{ref-busing_22}{2022}).

\section{Examples}\label{examples}

We analyze two classical MDS examples: the color data of Ekman
(\citeproc{ref-ekman_54}{1954}) and the Morse code data of Rothkopf
(\citeproc{ref-rothkopf_57}{1957}). We are only interested in this paper
in the properties of the programs, not in the substantive interpretation
of the results.

\subsection{Results}\label{results}

Numerical elastic scaling takes 586 iterations to arrive at stress
2.3268637. For ordinal we need 437 iterations for stress 0.056998. Since
the ordinal fit is very good the log-stress of \eqref{eq-approx} should
be close to the elastic stress. It is 0.0581521.

For the Morse code data numerical elastic scaling takes 212 iterations
to arrive at stress 64.3828862, for ordinal we need 152 iterations for
stress 29.7732277. Because the fit is much worse than in the Ekman case
we do not expect log-stress to be close to elastic stress. For the
ordinal case final log-stress is 42.8850382.

\subsection{Timing}\label{timing}

The results of a comparison of the R and C implementations with
microbenchmark (Mersmann (\citeproc{ref-mersmann_24}{2024})) yields the
times in microseconds in following table.

\begin{verbatim}
               min     lq   mean median     uq    max
R/Numerical  40.13  41.66  44.38  43.84  45.11 109.41
C/Numerical   1.44   1.53   1.60   1.60   1.64   1.81
R/Ordinal   163.52 168.14 172.10 169.50 172.28 242.27
C/Ordinal     4.06   4.39   4.57   4.46   4.71   7.37
\end{verbatim}

Using median times we see that that C version is about 40 times as fast
as the R version.

Next, the timings for the Morse data.

\begin{verbatim}
               min     lq   mean median     uq    max
R/Numerical  80.71  85.50  92.21  87.61  91.38 162.41
C/Numerical   6.38   6.62   6.77   6.71   6.80  10.96
R/Ordinal   711.81 744.80 767.76 764.32 784.45 862.41
C/Ordinal     7.25   7.60   7.82   7.74   7.83  11.35
\end{verbatim}

In the numerical case the C version is about 15 times as fast, in the
ordinal case about 100 times. The huge difference in the ordinal case is
likely due in large part to the different monotone regression routines
used by the R and C programs.

Of course these microbenchmark results depend on the default parameters
of the programs (same in R and C), on the speed of my computer (Mac Mini
with Apple M4 Pro, 64 GB of RAM, Tahoe 26.2, Apple clang 17.0.0), and on
my programming habits (same in R and C).

\subsection{Plots}\label{plots}

We use Shepard plots ({de Leeuw and Mair}
(\citeproc{ref-deleeuw_mair_C_15}{2015})) to illustrate the difference
between the solutions for classical Kruskal stress and elastic stress.
Both for Ekman and Morse we give the two solutions in the ordinal case,
with the primary approach to ties.

\begin{center}
\pandocbounded{\includegraphics[keepaspectratio]{elastic_files/figure-pdf/compplotE-1.pdf}}
\end{center}

\begin{center}
\pandocbounded{\includegraphics[keepaspectratio]{elastic_files/figure-pdf/compplotE-2.pdf}}
\end{center}

\begin{center}
\pandocbounded{\includegraphics[keepaspectratio]{elastic_files/figure-pdf/compplotM-1.pdf}}
\end{center}

\begin{center}
\pandocbounded{\includegraphics[keepaspectratio]{elastic_files/figure-pdf/compplotM-2.pdf}}
\end{center}

The examples show that the solutions are pretty similar, but definitely
elastic MDS gives a better fit for the smaller dissimilarities and
deemphasizes the larger ones. In that sense elastic MDS is more
``robust''.

\section*{References}\label{references}
\addcontentsline{toc}{section}{References}

\protect\phantomsection\label{refs}
\begin{CSLReferences}{1}{1}
\bibitem[\citeproctext]{ref-borg_groenen_05}
Borg, Ingwer, and Patrick J. F. Groenen. 2005. \emph{Modern
Multidimensional Scaling}. Second Edition. Springer.

\bibitem[\citeproctext]{ref-busing_22}
Busing, Frank M. T. A. 2022. {``{Monotone Regression: A Simple and Fast
O(n) PAVA Implementation.}''} \emph{Journal of Statistical Software} 102
(Code Snippet 1).
\url{https://www.jstatsoft.org/index.php/jss/article/view/v102c01/4306}.

\bibitem[\citeproctext]{ref-deleeuw_U_75a}
{de Leeuw, Jan}. 1975. {``{A Normalized Cone Regression Approach to
Alternating Least Squares Algorithms}.''} Department of Data Theory
FSW/RUL.
\url{https://deleeuw.github.io/janspubs/deleeuw-u-75-a/deleeuw-u-75-a.pdf}.

\bibitem[\citeproctext]{ref-deleeuw_E_25b}
{de Leeuw, Jan}. 2025. {``{Yet Another Smacof - Square Symmetric
Case}.''}
\url{https://deleeuw.github.io/janspubs/deleeuw-e-25-b/deleeuw-e-25-b.pdf}.

\bibitem[\citeproctext]{ref-deleeuw_hornik_mair_A_09}
De Leeuw, Jan, Kurt Hornik, and Patrick Mair. 2009. {``{Isotone
Optimization in R: Pool-Adjacent-Violators Algorithm (PAVA) and Active
Set Methods}.''} \emph{Journal of Statistical Software} 32 (5): 1--24.
\url{https://deleeuw.github.io/janspubs/deleeuw-hornik-mair-a-09/deleeuw-hornik-mair-a-09.pdf}.

\bibitem[\citeproctext]{ref-deleeuw_mair_A_09c}
De Leeuw, Jan, and Patrick Mair. 2009. {``{Multidimensional Scaling
Using Majorization: SMACOF in R}.''} \emph{Journal of Statistical
Software} 31 (3): 1--30.
\url{https://deleeuw.github.io/janspubs/deleeuw-mair-a-09-c/deleeuw-mair-a-09-c.pdf}.

\bibitem[\citeproctext]{ref-deleeuw_mair_C_15}
{de Leeuw, Jan, and Patrick Mair}. 2015. {``{Shepard Diagram}.''} In
\emph{Wiley {StatsRef}: Statistics Reference Online}. Wiley.
\url{https://deleeuw.github.io/janspubs/deleeuw-mair-c-15/deleeuw-mair-c-15.pdf}.

\bibitem[\citeproctext]{ref-ekman_54}
Ekman, Gosta. 1954. {``{Dimensions of Color Vision}.''} \emph{Journal of
Psychology} 38: 467--74.

\bibitem[\citeproctext]{ref-mair_groenen_deleeuw_A_22}
{Mair, Patrick, Patrick J. F. Groenen, and Jan de Leeuw}. 2022. {``{More
on Multidimensional Scaling in R: smacof Version 2}.''} \emph{Journal of
Statistical Software} 102 (10): 1--47.
\url{https://deleeuw.github.io/janspubs/mair-groenen-deleeuw-a-22/mair-groenen-deleeuw-a-22.pdf}.

\bibitem[\citeproctext]{ref-mcgee_65}
McGee, Victor E. 1965. {``More on an "Elastic" Multidimensional Scaling
Procedure.''} \emph{Perceptual and Motor Skills} 21: 81--82.

\bibitem[\citeproctext]{ref-mcgee_66}
McGee, Victor E. 1966. {``The Multidimensional Analysis of 'Elastic'
Distances.''} \emph{British Journal of Mathematical and Statistical
Psychology} 19 (2): 181--96.

\bibitem[\citeproctext]{ref-mcgee_67}
McGee, Victor E. 1967. {``A Reply to Some Criticisms of Elastic
Multidimensional Scaling.''} \emph{British Journal of Mathematical and
Statistical Psychology} 20 (2): 243--47.

\bibitem[\citeproctext]{ref-mcgee_68}
McGee, Victor E. 1968. {``{Multidimensional Scaling of N Sets of
Similarity Measures: A Nonmetric Individual Differences Approach}.''}
\emph{Multivariate Behavioral Research} 3 (2): 233--48.

\bibitem[\citeproctext]{ref-mersmann_24}
Mersmann, Olaf. 2024. \emph{{microbenchmark: Accurate Timing
Functions}}. \url{https://CRAN.R-project.org/package=microbenchmark}.

\bibitem[\citeproctext]{ref-ramsay_77}
Ramsay, James O. 1977. {``{Maximum Likelihood Estimation in
Multidimensional Scaling}.''} \emph{Psychometrika} 42: 241--66.

\bibitem[\citeproctext]{ref-rothkopf_57}
Rothkopf, Ernst Z. 1957. {``{A Measure of Stimulus Similarity and Errors
in some Paired-associate Learning}.''} \emph{Journal of Experimental
Psychology} 53: 94--101.

\bibitem[\citeproctext]{ref-rusch_deleeuw_chen_mair_M_25}
{Rusch, Thomas, Jan de Leeuw, Lisha Chen, and Patrick Mair}. 2025.
\emph{{Package "smacofx": Flexible Multidimensional Scaling and 'smacof'
Extensions}}.
\url{https://deleeuw.github.io/janspubs/rusch-deleeuw-chen-mair-m-25/rusch-deleeuw-chen-mair-m-25.pdf}.

\bibitem[\citeproctext]{ref-rusch_deleeuw_mair_hornik_26}
{Rusch, Thomas, Jan de Leeuw, Patrick Mair, and Kurt Hornik}. 2026.
{``{Flexible Multidimensional Scaling with the R Package smacofx}.''}
\emph{Journal of Statistical Software}.
\url{https://deleeuw.github.io/janspubs/rusch-deleeuw-mair-hornik-a-25/rusch-deleeuw-mair-hornik-a-25.pdf}.

\bibitem[\citeproctext]{ref-torgerson_58}
Torgerson, Warren S. 1958. \emph{{Theory and Methods of Scaling}}.
Wiley.

\end{CSLReferences}




\end{document}
